\documentclass[letterpaper]{article}

% Use the official AAAI style file when preparing the camera-ready package.
% \usepackage{aaai27}
\usepackage{times}
\usepackage{helvet}
\usepackage{courier}
\usepackage{amsmath}
\usepackage{amssymb}
\usepackage{booktabs}
\usepackage{graphicx}
\usepackage{multirow}
\usepackage{xcolor}
\usepackage{algorithm}
\usepackage{algorithmic}

\title{Support Density Matters: Density-Guided Feature Tracing for Training-Free Open-Vocabulary Semantic Segmentation}

\author{
Anonymous Authors
}

\begin{document}
\maketitle

\begin{abstract}
Open-vocabulary semantic segmentation requires dense visual predictions from
vision-language models that are primarily trained with image-level supervision.
Recent training-free pipelines improve this transfer by reusing frozen CLIP
features and avoiding dataset-specific optimization. However, we observe that
the decisive factor for dense transfer is not only whether a patch feature is
semantically aligned, but also whether it is sufficiently supported by nearby
visual evidence on the patch manifold. Low-support patches frequently appear
near object boundaries, thin structures, and small regions, where direct
text-image similarity produces fragmented predictions. We therefore propose
DeFT, a Density-guided Feature Tracing framework for training-free
open-vocabulary segmentation. DeFT estimates support density among frozen visual
tokens, traces reliable high-density evidence to low-support regions, and
performs structure-conserving multi-scale recalibration before text-image
similarity rendering. The whole procedure is deterministic and introduces no
learnable parameters. DeFT can be inserted into existing CLIP/SFP-style
backbones without retraining. Experiments on standard open-vocabulary
segmentation benchmarks show consistent improvements over the backbone,
especially on boundary-sensitive and small-object categories, while preserving
the training-free property.
\end{abstract}

\section{Introduction}

Open-vocabulary semantic segmentation (OVSS) aims to segment arbitrary
categories described by natural language. Vision-language models such as CLIP
provide strong image-text alignment, but their patch features are not directly
optimized for pixel-level prediction. A common training-free solution is to
freeze the vision-language model, extract patch-level visual features, and
compute dense similarity with text embeddings. This design is appealing because
it avoids dataset-specific training and preserves broad vocabulary coverage.

Despite this appeal, dense transfer from CLIP remains structurally fragile. In
many images, confident semantic responses concentrate on a small number of
patches, while boundaries, thin parts, and small objects receive weak support
from their surrounding token neighborhood. We call this phenomenon \emph{support
density imbalance}. A patch can be semantically meaningful but still unreliable
for dense prediction when its local and non-local neighbors provide insufficient
evidence. This perspective shifts the focus from isolated token correction to
the density of visual evidence available for each region.

In this paper, we ask a simple question: can a frozen segmentation pipeline be
improved by explicitly tracing reliable support from dense regions to sparse
ones? We answer yes and propose \textbf{DeFT}, a \textbf{De}nsity-guided
\textbf{F}eature \textbf{T}racing framework. Given frozen visual patch features,
DeFT constructs a token support graph using cosine similarity, estimates a
support density score for every patch, and computes a density gate that
identifies regions requiring structural assistance. It then aggregates evidence
from semantically similar tokens and spatial neighborhoods, followed by a
multi-scale recalibration step that restores local continuity without learning
new parameters.

DeFT is designed as a plug-in for existing training-free OVSS backbones. In our
implementation, we insert DeFT after the backbone produces purified visual
features and before text-image similarity is computed. This placement preserves
the original semantic space while improving spatial coherence. Since DeFT has no
learnable weights and uses fixed hyperparameters, it keeps the training-free
protocol intact.

Our contributions are summarized as follows:
\begin{itemize}
    \item We identify support density imbalance as a key obstacle in
    training-free open-vocabulary semantic segmentation, especially for
    boundaries, thin structures, and small regions.
    \item We propose DeFT, a deterministic feature tracing framework that
    estimates patch support density and propagates reliable evidence to
    low-support regions without training.
    \item We introduce a structure-conserving multi-scale recalibration step
    that improves dense prediction while preserving CLIP's original semantic
    embedding space.
    \item We provide comprehensive comparisons, ablations, and visual analyses
    to verify accuracy, boundary quality, small-object behavior, and runtime
    cost.
\end{itemize}

\section{Related Work}

\paragraph{Open-vocabulary semantic segmentation.}
OVSS extends semantic segmentation to categories unseen during segmentation
training. Early methods adapt vision-language models with additional supervision,
while recent approaches explore training-free inference by directly reusing
frozen image-text representations. DeFT follows the training-free line and does
not introduce dataset-specific optimization.

\paragraph{Dense transfer from vision-language models.}
CLIP-like models learn image-level alignment and contain transferable visual
tokens, yet their features are not guaranteed to be spatially coherent. Existing
methods improve dense transfer through prompt design, feature refinement,
attention manipulation, or post-processing. DeFT is complementary: it operates
on the frozen visual tokens and explicitly models how much support each patch
receives from the visual token manifold.

\paragraph{Graph and spatial refinement.}
Graph-based propagation and spatial smoothing are widely used in dense
prediction. Unlike trainable graph neural networks or learned refinement heads,
DeFT uses a fixed similarity graph and deterministic multi-scale operations.
This makes it suitable for training-free OVSS, where introducing learnable
modules would weaken the protocol.

\section{Method}

Let $X \in \mathbb{R}^{N \times C}$ denote frozen visual patch features produced
by a CLIP/SFP-style backbone, where $N=H \times W$. Let $T \in
\mathbb{R}^{K \times C}$ denote normalized text embeddings for $K$ candidate
classes. The baseline computes dense logits by cosine similarity between $X$
and $T$. DeFT refines $X$ before this similarity step:
\begin{equation}
    \hat{X} = \mathrm{DeFT}(X), \quad
    Y = \hat{X} T^\top.
\end{equation}
All operations in DeFT are fixed and contain no learnable parameters.

\subsection{Support Density Estimation}

We first build a patch support graph in the normalized visual feature space. For
each token $x_i$, we compute cosine similarity to all other tokens:
\begin{equation}
    s_{ij} = \frac{x_i^\top x_j}{\|x_i\|_2 \|x_j\|_2}.
\end{equation}
Let $\mathcal{N}_k(i)$ be the top-$k$ most similar neighbors of token $i$. We
define the support density of $x_i$ as:
\begin{equation}
    d_i = \frac{1}{k} \sum_{j \in \mathcal{N}_k(i)} \max(s_{ij}, 0).
\end{equation}
The density is min-max normalized per image. A density gate is then obtained by:
\begin{equation}
    g_i = \sigma \left( \frac{\tau - d_i}{\gamma} \right),
\end{equation}
where $\tau$ is the density threshold and $\gamma$ is the temperature. A larger
$g_i$ indicates that token $i$ has lower visual support and should receive more
traced evidence.

\subsection{Density-Guided Evidence Tracing}

For each token, DeFT gathers two sources of support. The first source is a
semantic graph aggregate:
\begin{equation}
    z_i^{g} = \sum_{j \in \mathcal{N}_k(i)}
    \frac{\exp(s_{ij}/\beta)}{\sum_{m \in \mathcal{N}_k(i)} \exp(s_{im}/\beta)}
    x_j,
\end{equation}
where $\beta$ controls graph sharpness. The second source is a local spatial
aggregate $z_i^{l}$ computed by average pooling on the reshaped feature map over
multiple window sizes. The traced context is:
\begin{equation}
    z_i = \lambda_g z_i^g + \lambda_l z_i^l.
\end{equation}
The output token is a density-gated mixture:
\begin{equation}
    x_i' = (1 - \alpha g_i) x_i + \alpha g_i z_i.
\end{equation}
Thus, high-density tokens remain close to the original backbone representation,
while low-density tokens borrow evidence from reliable semantic and local
neighbors.

\subsection{Structure-Conserving Multi-Scale Recalibration}

Evidence tracing improves token reliability but may still leave local
discontinuities. DeFT therefore applies a deterministic multi-scale
recalibration to the feature map $X'$. For kernel sizes $\mathcal{K}$, we compute
multi-scale residuals:
\begin{equation}
    r = \frac{1}{|\mathcal{K}|} \sum_{q \in \mathcal{K}}
    \left(\mathrm{AvgPool}_q(X') - X'\right).
\end{equation}
The residual is injected more strongly into low-density regions:
\begin{equation}
    \tilde{X} = X' + \eta (1-D) \odot r,
\end{equation}
where $D$ is the density map and $\eta$ is a fixed strength. Finally, features
are $\ell_2$-normalized to preserve the cosine similarity interface with text
embeddings.

\subsection{Text-Conditioned Confidence Rendering}

The refined features are compared with text embeddings using the original
backbone classifier. Optionally, DeFT renders logits with a mild support
confidence prior:
\begin{equation}
    Y_{k,i} = (\tilde{x}_i^\top t_k) \cdot
    \left(1 + \rho(d_i - \bar{d})\right),
\end{equation}
where $\rho$ is fixed and small. This step does not change the class vocabulary
or text prompts. In the main experiments, we report both with and without this
rendering term to isolate its effect.

\begin{algorithm}[t]
\caption{DeFT Inference}
\begin{algorithmic}[1]
\STATE Input frozen visual tokens $X$, text embeddings $T$
\STATE Normalize visual tokens and build cosine similarity graph
\STATE Estimate support density $D$ from top-$k$ positive neighbors
\STATE Compute density gate $G$ from $D$
\STATE Aggregate semantic graph evidence $Z^g$
\STATE Aggregate local multi-scale evidence $Z^l$
\STATE Trace evidence to low-density tokens using $G$
\STATE Apply structure-conserving multi-scale recalibration
\STATE Compute logits with original text embeddings
\STATE Output segmentation map
\end{algorithmic}
\end{algorithm}

\section{Experiments}

\subsection{Experimental Setup}

\paragraph{Datasets.}
We evaluate on standard OVSS benchmarks, including PASCAL VOC, PASCAL Context,
COCO-Stuff, and ADE20K. Dataset-specific details follow the protocol of the
backbone method.

\paragraph{Backbone.}
We use the official SFP-style training-free backbone as the base pipeline. DeFT
is inserted after visual feature purification and before text-image similarity.
All backbone settings, prompts, and post-processing are kept unchanged unless
explicitly stated.

\paragraph{Metrics.}
We report mean Intersection-over-Union (mIoU), mean accuracy (mAcc), boundary
F-score, small-object mIoU, runtime, and GPU memory. Boundary and small-object
metrics are used to verify the structural claim rather than only reporting
global mIoU.

\paragraph{Implementation details.}
Unless otherwise specified, DeFT uses $k=12$, $\tau=0.46$, $\alpha=0.72$,
$\lambda_g=0.65$, $\lambda_l=0.35$, and multi-scale kernels $\{3,5,7\}$.
The same fixed hyperparameters are used across datasets.

\subsection{Main Results}

\begin{table*}[t]
\centering
\caption{Main comparison on open-vocabulary semantic segmentation benchmarks.
Fill the empty cells after running the provided evaluation matrix.}
\begin{tabular}{lcccccc}
\toprule
Method & Training-Free & PASCAL VOC & PASCAL Context & COCO-Stuff & ADE20K & Avg. \\
\midrule
CLIP baseline & Yes & -- & -- & -- & -- & -- \\
MaskCLIP-style baseline & Yes & -- & -- & -- & -- & -- \\
SFP backbone & Yes & -- & -- & -- & -- & -- \\
\textbf{DeFT (ours)} & Yes & -- & -- & -- & -- & -- \\
\bottomrule
\end{tabular}
\end{table*}

\subsection{Ablation Study}

\begin{table}[t]
\centering
\caption{Ablation of DeFT components.}
\begin{tabular}{lccc}
\toprule
Variant & mIoU & Boundary-F & Small mIoU \\
\midrule
SFP backbone & -- & -- & -- \\
DeFT w/o graph evidence & -- & -- & -- \\
DeFT w/o local evidence & -- & -- & -- \\
DeFT w/o multi-scale recalibration & -- & -- & -- \\
DeFT full & -- & -- & -- \\
\bottomrule
\end{tabular}
\end{table}

\begin{table}[t]
\centering
\caption{Sensitivity to the number of support neighbors $k$.}
\begin{tabular}{lccccc}
\toprule
$k$ & 4 & 8 & 12 & 16 & 24 \\
\midrule
mIoU & -- & -- & -- & -- & -- \\
Boundary-F & -- & -- & -- & -- & -- \\
\bottomrule
\end{tabular}
\end{table}

\subsection{Qualitative Analysis}

We visualize support density maps, tracing gates, and segmentation maps before
and after DeFT. The expected pattern is that low-density gates concentrate on
ambiguous boundaries, thin structures, and small regions. After tracing,
activation maps should become more spatially continuous while preserving class
specificity. We also visualize token similarity graphs to show how DeFT traces
evidence from high-support regions to sparse regions.

\section{Conclusion}

We presented DeFT, a training-free framework for open-vocabulary semantic
segmentation. Instead of adding a trainable decoder, DeFT studies the support
density of frozen visual tokens and traces reliable evidence to regions with
insufficient support. By combining graph evidence, local evidence, and
multi-scale structure recalibration, DeFT improves dense prediction while
preserving the original vision-language embedding space. Future work can explore
faster approximate support graphs and stronger prompt-conditioned rendering.

\bibliographystyle{aaai}
\bibliography{references}

\end{document}
